% !TeX root = ../papers/nc-p3.tex
\section{被約二十点の五空間における階数一曲線と反対対称性}
\label{sec:finite-five-space-rank-one-curves}
\begin{theorem}\label{thm:finite-five-space-rank-one-curves}
定理\ref{thm:finite-marked-five-space-secants}の被約二十点の五空間では、
少なくとも一方の反対位置にテンソル対称性がある。
従って被約な有限窓で全中央支持が4、反対階数10なら片側対称であり、反対階数11は生じない。
\end{theorem}
\begin{proof}[母線の排除]
前節の記号を使う。$E_A,E_B$ はそれぞれ非零の片側評価が全て階数2以上で、
共通の $A(p)=B(p)$ は可逆である。後者は印での行列の階数が
$3-\dim T_pX$ であることによる。左右を交換して $E$ に母線 $ku\otimes T$、
$\dim T=2$ があると仮定する。行列の像を $ku$ で割った二行三列空間
$\bar E$ の次元は2または3である（像が $ku$ に入る行列は高々三次元）。
前節の独立な鉛筆を与える13点以上のうち、左核 $l$ が $l(u)=0$ を
満たす点は印を通る位置1の平面にあり、印を除いて高々6点である。
残る各点の鉛筆は $\bar E$ に単射し、射影後も共通右核 $r$ がある。
$\dim\bar E=2$ なら鉛筆の像は全体で、その中に階数2の行列があるため $r$ は一意。
固定 $r$ は印からの射影の一方向なので位置3で高々2点しか与えない。

$\dim\bar E=3$ とする。射影平面の階数一軌跡が有限なら、それは
$\mathbb P^1\times\mathbb P^2$ の三超平面完全交叉で長さ3である。
各接空間は第一因子への射影で単射である（核に非零方向があれば母線が丸ごと入る）。
従って曲線状で長さ2の部分スキームは高々三つ。
共通右核を持つ鉛筆は二行二列の行列式でそのような部分を切り、その線型包に等しい。
各々の右核が一意なので高々6点となり、残る7点以上に足りない。
正次元なら整曲線の両射影の包のベクトル次元 $a,b$ に対し
Segre像の包の次元は $a+b-1$ 以上である
（定理\ref{thm:finite-marked-ruling-pencil}の付値による積次元評価）。
第一因子が固定された母線はその消滅行による $E$ の評価階数を1以下にする。
平面全体が一つの固定因子を持つことも階数2の元に反する。
両因子とも動くなら $a=b=2$ で曲線が平面を張り、全 $\bar E$ が固定した
二列空間に入る。その消滅列での $E$ の評価像は $ku$ 以下となって矛盾する。
従って第二因子の固定された母線 $k^2\otimes v$ が入る。
$r(v)\ne0$ なら $\bar E$ の $r$ 評価は全射で核は一次元にすぎない。
全ての独立な鉛筆の点が $l(u)=0$ または $r(v)=0$ に入り、
それぞれ印以外に高々6点なので再び13点には足りない。
\end{proof}
\begin{proof}[平面二次曲線の場合]
階数一軌跡には前節より整曲線がある。それが全体を張れば
命題\ref{prop:single-quartic-slice-kernel-symmetry}を適用する。
それ以外の曲線では両因子が動く。実際、固定因子の線型空間と $\mathbf P(E)$ の交叉は線型であり、
曲線を含めば母線を含む。両射影が平面を張れば積次元評価により全体を張る。
一方が直線を張ればその法線での評価核は三次元以下であり、
同じ評価から他方も直線を張る。このとき曲線は二行二列のSegre二次曲面の平面切断で、
母線がないから滑らかな二次曲線である。

この曲線の包を $D$ とし、左右の基底を選んで $D$ を左上二行二列の対称行列全体とする。
$E/D$ の同次座標を $y=(s,t)$ として一つの線型補空間を選び、その行列を
\[
 M(y)=\begin{pmatrix}K(y)&c(y)\\ d(y)^t&h(y)\end{pmatrix}
\]
とする。$y\mapsto (c,h)$ と $y\mapsto(d,h)$ はともに単射である。最後の列または行で
$E$ を評価した階数が2以上だからである。$[N]$ が $D$ の外の階数一点で $y$ がその商座標なら、
$h(y)\ne0$ である。零なら最終列と最終行がともに非零なのに隅成分が零で階数一に反する。
左上ブロックは $c d^t/h$ に決まる。よってこの点の存在条件は二次式
\[
 Q(y)=c_1d_2-c_2d_1-h(K_{12}-K_{21})=0
\]
であり、点も $[y]$ から一意である。

$Q\ne0$ とする。二つまでの根から次のように高々二点を選ぶ。
$h(y)\ne0$ の根からは上の階数一点を選ぶ。$h(y)=0$ の根では
非零な $c(y),d(y)$ が比例するので二次曲線上の $[c(y)c(y)^t]$ を選ぶ。
この選択は曲線上の接線で $D$ の外へ出るものも捉える。
実際 $[vv^t]$ での階数一行列の接方向は、一方の因子だけを動かす二項の和である。
その最終行・列は $h=0$ かつ $c,d\in kv$ を要求するため、商座標が非零なら上の根から得た点である。

各 $q$ で $A(q),B(q)$ が比例すれば、その任意の共通商への像も比例する。
独立なら鉛筆は共通左右核を持ち、母線がないので行列式はその上で長さ2の階数一因子を切る。
その因子に $D$ 外の点があれば鉛筆は選んだ点を通る。相異なる二点が $D$ 内にあれば鉛筆も
$D$ 内にある。一点に集中すれば接線であり、$D$ 外に出る場合も選んだ点を通る。
一方 $D$ 内の独立な鉛筆の共通左右核は二行二列部分の外の固定方向である。
実際二次元対称行列の空間で指定した非零二列を殺す部分は一次元である。
従ってこの例外点は位置1で印を通る固定直線にあり、印自身を除いて高々2点である。

選んだ点の包を含む任意のベクトル二平面 $T\subset E$ を取り $F=E/T$ とする。
$A,B$ の $F$ への像の外積三成分は例外の高々2点だけに支持される。
これらは形式的な位置2と0の双一次式としては三つ独立である。
実際 $F_A,F_B\subset F$ を両写像の像とすれば各次元は2以上で和は三次元。
交代形式が $F_A\times F_B$ 上で消えるとする。両方が平面なら異なる平面で、
共通の基底元を含む適合基底の三対を全て消す。一方が全空間なら根基が二次元以上になり、
やはり零である。よって二十点への双一次評価は単射でなく、
定理\ref{thm:finite-opposite-product-gap}が偶位置の対称性を与える。

$Q=0$ が恒等的で $h\ne0$ の場合、$y$ を動かした
$(c(y),h(y))(d(y),h(y))^t$ はもう一つの滑らかな二次曲線となる。
実際両列写像は単射であり、二次Veroneseを与える。この曲線を張る平面と最初の平面は
ちょうど一点で交わる。二つの行支持および列支持はいずれも異なる二次元空間であり、
共通部分が入る一行一列空間の次元は1、二平面は射影四空間で必ず交わるためである。
この二曲線はその一点で交わる（共通行列は階数一で、第二の曲線については
$h=0$ の値を取る）。その被約和 $C$ は連結で $\mathbf P(E)$ を張る。
$E_A$ はどちらの曲線も含めない（単独の四空間では正次元切断が母線だけである）。
従って超平面切断は非零因子であり、$H^0(C,\mathcal O_C)=k$ からその切断は
$\mathbf P(E_A)$ を張る。長さ4の判定は前節の命題の証明と同じで核側対称性が従う。

最後に $Q=h=0$ がともに恒等式なら $c,d$ は二次元への同型で、$d$ は $c$ の定数倍である。
最後の行の倍率を揃え、さらに最後の行の定数倍を上二行に足せば $K_{12}-K_{21}$ を消せる。
実際 $c$ の二成分は独立であり、この操作で任意の一次の歪成分を加えられる。
こうして $E$ 全体が対称となる。$D$ 外の階数一点はない。
$E_A$ は $D$ を含めないから、その階数一点は二次曲線上の高々二点である。
$A(q)$ が階数2なら射影した左右核は一致し、印では両射影ベクトルは零である。
従って左右三座標の外積も高々二点に支持され、三つの独立な奇位置の双一次式から
同じ積次元定理で核側対称性が従う。

反対階数については、印の両側の和の次元の和は反対階数以下である
（定理\ref{thm:finite-rank-nine-all-marks}の直交性の証明）。
従って階数11以下なら一方の和は4または5である。4なら同定理の曲線状印での
共通切断の議論、5なら本定理を使う。片側対称なら反対切断は四次元の対称二乗を経由するため階数は10以下である。
\end{proof}
